% Chapter 3
\chapter{System Architecture and Application Design}
\label{Chapter3}
\lhead{Chapter 3. \emph{System Architecture and Application Design}}

The software developed in this research forms the core of the Magnavis application. Built as a desktop application using Python and PyQt5, it enables interactive handling of magnetic field data, prediction, and anomaly detection. This chapter describes the system design, main modules, and how the GUI integrates multiple functionalities.

\section{High-Level Architecture}
Magnavis comprises two main entry points:
\begin{enumerate}[label=(\alph*)]
    \item \textbf{application.py} --- Legacy USGS real-time workflow.
    \item \textbf{app.py} --- Primary DB/CSV workflow used in benchmarks: multi-sensor support, pretrained GRU/LSTM loading, and headless batch mode for evaluation scripts.
\end{enumerate}

Both variants share the base classes \texttt{Application} and \texttt{ApplicationWindow} from \texttt{application.py}; the temp variant replaces the USGS fetch with \texttt{data\_convert\_db\_now} and extends the UI with a two-column layout and direction finding.

\section{Main Modules and Responsibilities}

\subsection{application.py (Orchestrator)}
\begin{itemize}
    \item Builds and runs the Qt application and main window.
    \item Loads UI files (\texttt{ui\_ApplicationWindow.ui}, \texttt{ui\_MagneticTimeSeriesWidget.ui}).
    \item Manages data fetch threads (QThread + worker), plot setup and periodic updates (Matplotlib timers), VTK pipeline for spatial visualisation, session folder creation, prediction subprocess lifecycle, and anomaly detection calls.
\end{itemize}

\subsection{data\_convert\_now.py (Real-Time Ingestion)}
\begin{itemize}
    \item Downloads magnetic field data from USGS GeoMag (element \texttt{H}, station \texttt{BRW} in the default configuration).
    \item Writes raw JSON to the session folder as \texttt{download\_mag.json}.
    \item Returns a DataFrame with columns \texttt{time\_H} and \texttt{mag\_H\_nT}.
\end{itemize}

\subsection{data\_convert\_db\_now.py (DB-Backed Multi-Sensor)}
\begin{itemize}
    \item Fetches time-series from the project MySQL database (or CSV/Simulation) for selected sensors.
    \item Provides \texttt{fetch\_timeseries\_window\_multi}, \texttt{get\_timeseries\_magnetic\_data\_multi}, and related functions used by \texttt{application\_temp.py}.
\end{itemize}

\subsection{predictor\_ai.py (Prediction Subprocess)}
\begin{itemize}
    \item Standalone script invoked as \texttt{python src/predictor\_ai.py <path/to/predict\_input.csv>}.
    \item Reads \texttt{predict\_input.csv} (columns \texttt{x}=timestamps, \texttt{y}=values).
    \item Uses a GRU-based Keras model (GRUPredictor) with cyclic time-of-day features; optionally loads a pre-trained model via \texttt{PRETRAINED\_MODEL\_PATH}.
    \item Writes \texttt{predict\_out.csv} to the session folder.
    \item Respects \texttt{TRAIN\_WINDOW\_MINUTES} to restrict training to the most recent N minutes.
\end{itemize}

\subsection{Anomaly\_detector.py (Anomaly Detection)}
\begin{itemize}
    \item Class \texttt{AnomalyDetector}: interpolates predicted values at actual timestamps, maintains prediction error history, computes adaptive threshold (mean + multiplier $\times$ std), and supports an optional freeze window after the first anomaly.
    \item Returns anomalies and current threshold; used by the main application to flag and visualise anomalies.
\end{itemize}

\subsection{anomaly\_direction.py (Direction Finding)}
\begin{itemize}
    \item \texttt{compute\_direction\_obs1()} / \texttt{compute\_direction\_obs2()}: compute azimuth and inclination from component perturbations (baseline subtracted).
    \item \texttt{triangulate\_source\_location()}: when both observatories have a recent anomaly, estimates approximate source position from bearing lines.
\end{itemize}

\section{UI Layout (application\_temp.py)}
\begin{itemize}
    \item \textbf{Left half:} Up to three sensor stream panels, each in a QScrollArea, showing time-series of B (nT) --- low-pass filtered resultant magnitude.
    \item \textbf{Right half:} (1) Shared parameters panel (threshold multiplier, training window, freeze window); (2) 3D anomaly direction plot (unit vectors from observatory centre); (3) Log (scrollable).
\end{itemize}

\section{Session Artifacts}
Under \texttt{src/sessions/<session\_id>/} (and per sensor in the temp variant):
\begin{itemize}
    \item \texttt{download\_mag.json} (real-time mode).
    \item \texttt{predict\_input.csv}, \texttt{predict\_out.csv} (anomalous timestamps excluded from input when possible).
    \item \texttt{predict\_stdout.log}, \texttt{predict\_stderr.log}.
\end{itemize}

\section*{Summary}
This chapter presented the architecture of Magnavis: the orchestrator (\texttt{application.py}), data ingestion (USGS or DB), GRU predictor subprocess, anomaly detector, and direction-finding module. The DB-backed variant extends the UI to multiple sensors and 3D direction visualisation while reusing the same prediction and anomaly logic.
